\chapter{От требований физики к планковскому миру}
\label{ch:descent}

Закон, записанный мельче уже проверенного масштаба, не свободен.
После осреднения он обязан вернуть те уравнения и те запреты, которые прибор уже сверил.
Ниже мы спускаемся по этажам и на каждом записываем, во что это обязательство превращается.
Планковский мир в конце есть то место, где непрерывное поле в точке перестаёт с обязательством справляться.

\section{Требования физики}
\label{sec:requirements}

Обязательство одно: нижний этаж воспроизводит верхний.
Чтобы было с чем сверять, скажем, что на макроскопическом масштабе установлено прямо, без гипотезы об устройстве вещества.

Возмущение доходит за конечное время.
Если бы влияние было мгновенным, показание прибора в данном месте зависело бы от того,
что в тот же момент происходит сколь угодно далеко, и локальный опыт не сходился бы.
За время~$\Delta t$ сигнал проходит путь не длиннее~$c\Delta t$:
\begin{equation}
  |\Delta\mathbf{x}| \le c\,\Delta t.
  \label{eq:light-cone}
\end{equation}
Любой более глубокий закон обязан сохранить этот конус.

Из \eqref{eq:light-cone} следует локальность.
Далёкое не успевает прийти, поэтому новое состояние в данной точке собирается из поля в её окрестности.
Для электромагнитного поля это уравнения Максвелла:
производная по времени стоит в точке, ротор и дивергенция берут соседние значения,
\begin{align}
  \nabla\times\mathbf{E} &= -\partial_t\mathbf{B}, &
  \nabla\cdot\mathbf{B} &= 0, \label{eq:maxwell-hom}\\
  \nabla\times\mathbf{B}
    &= \mu_0\mathbf{j} + \mu_0\varepsilon_0\,\partial_t\mathbf{E}, &
  \nabla\cdot\mathbf{E} &= \rho/\varepsilon_0. \label{eq:maxwell-inh}
\end{align}
Характеристики этой системы лежат на конусе \eqref{eq:light-cone}.
Та же схема у гидродинамики.
Для идеальной жидкости скорость в точке меняется от давления и скорости рядом:
\begin{equation}
  \partial_t\mathbf{v} + (\mathbf{v}\cdot\nabla)\mathbf{v}
  = -\frac{1}{\rho}\nabla p.
  \label{eq:euler}
\end{equation}
В механике сплошной среды справа стоит дивергенция напряжений,
$\partial_k\sigma_{ik}$, снова из окрестности той же точки.

Механическое состояние исчерпывается координатами и скоростями.
Положение точки есть $\mathbf{r}(t)$, скорость --- его производная,
\begin{equation}
  \mathbf{v} = \dot{\mathbf{r}}.
  \label{eq:velocity}
\end{equation}
Уравнения Ньютона связывают ускорение с этим состоянием:
\begin{equation}
  m\ddot{\mathbf{r}} = \mathbf{F}(\mathbf{r},\dot{\mathbf{r}}).
  \label{eq:newton}
\end{equation}
Это уравнение второго порядка.
Задание $\mathbf{r}$ и~$\dot{\mathbf{r}}$ в один момент восстанавливает их и в предыдущий, и в следующий.
Старшей производной среди начальных данных в \eqref{eq:newton} нет:
закон, которому она понадобилась бы, описывал бы другой мир.

Тепловое равновесие добавляет два запрета, тоже измеренных.
В адиабатически замкнутой системе энтропия не убывает,
\begin{equation}
  \mathrm{d}S \ge 0.
  \label{eq:second-law}
\end{equation}
При $T\to 0$ энтропия равновесной системы перестаёт зависеть от параметров равновесия~$\lambda$
и в шкале Планка обращается в нуль:
\begin{equation}
  \Bigl(\frac{\partial S}{\partial\lambda}\Bigr)_{T}
    \xrightarrow[T\to 0]{} 0,
  \qquad
  S \xrightarrow[T\to 0]{} 0.
  \label{eq:third-law}
\end{equation}
Система садится в одно основное состояние.
Закон, у которого при нулевой температуре оставалось бы макроскопически много
равноправных состояний, \eqref{eq:third-law} противоречит.

Вакуум, насколько его видит свет, не выделяет направления:
скорость одна и та же вдоль любой оси.
Вместе с \eqref{eq:light-cone} это интервал
\begin{equation}
  \mathrm{d}s^2 = c^2\,\mathrm{d}t^2 - \mathrm{d}\mathbf{x}^2,
  \label{eq:interval}
\end{equation}
один и тот же во всех инерциальных системах.
Заряд в замкнутой системе сохраняется.
Это уравнение непрерывности
\begin{equation}
  \partial_t\rho + \nabla\cdot\mathbf{j} = 0,
  \label{eq:continuity}
\end{equation}
а на узле проводников --- закон Кирхгофа, $\sum_k I_k = 0$.
Не свойство, приписанное заранее микроскопическому полю.

На этом макроскопическая сверка кончается.
Спин, комплексная амплитуда, целочисленность устойчивого заряда
и вопрос о том, откуда берутся безразмерные числа, здесь ещё не видны.
Они появятся на том этаже, где без них уже проверенный опыт перестаёт получаться,
и каждое окажется тем же конусом, тем же сохранением или той же обратимостью, дочитанными до этого этажа.

\section{Макромир}
\label{sec:macro-known}

Макромир есть эти же уравнения, записанные для непрерывных полей.

Из \eqref{eq:second-law} и \eqref{eq:third-law} строится термодинамика.
Температура есть производная энергии по энтропии,
\begin{equation}
  T = \Bigl(\frac{\partial E}{\partial S}\Bigr)_{V,N},
  \label{eq:temperature}
\end{equation}
и из неё же --- потенциалы.
Интервал \eqref{eq:interval} один и тот же во всех инерциальных системах,
поэтому энергия и импульс составляют четырёхвектор
\begin{equation}
  p^{\mu} = (E/c,\,\mathbf{p}).
  \label{eq:four-momentum}
\end{equation}

Из этих уравнений собрана система единиц.
Числа~$c$, $\hbar$, $k_B$, заряд~$e$, массы частиц и постоянная тонкой структуры~$\alpha_{\mathrm{fs}}$ в ней измерены.
Сами уравнения их не выводят.
На этом этаже такое ещё терпимо: число может быть константой прибора.
Один этаж ниже тот же закон, если он претендует быть фундаментальным, уже не имеет права оставлять число вписанным рукой.

\section{Микромир}
\label{sec:macro-open}
\label{sec:microworld}

Внутри вещества непрерывные уравнения, только что записанные, ещё держат механику и электродинамику прибора.
Они перестают держать сам опыт, как только отсчёт становится единичным.

Ансамбль подготовленных одинаково атомов даёт частоты исходов.
Один запуск прибора даёт один исход.
Уравнения поля содержат частоты и не содержат границу между системой, прибором и записью результата.
Сохранение того, что в ансамбле выглядит как полная вероятность, здесь ещё не объясняет одиночный отсчёт.

Числа~$c$, $\hbar$, $G$, $e$, $\alpha_{\mathrm{fs}}$ и массы фермионов измерены с большой точностью.
Безразмерное число, стоящее в законе свободным множителем, есть выбор, а не следствие.
Пока закон претендует только описать прибор, такой множитель терпим.
Как только тот же закон считается фундаментальным, вписанное число означает, что измерение не выведено.

Чёрная дыра делает ограниченность плотности видимой числом.
Энтропия горизонта пропорциональна его площади, а не объёму под ним.
Число различимых состояний в причинно связанной области поэтому конечно:
оно не может расти так, как рос бы набор независимых значений поля в каждой точке объёма.
Энергия в той же области ограничена сверху по той же причине:
слишком большая энергия замыкает область собственным горизонтом.
Общая теория относительности и квантовая теория поля локально друг другу не противоречат и этого счёта не содержат.

Там, где одновременно существенны $\hbar$, $G$ и~$c$,
поле в каждой точке этот счёт нарушает.
Следующий этаж --- запись того же сохранения и той же обратимости,
при которой спектр атома и два знака спина уже получаются.

\section{Квантовая механика}
\label{sec:qm}

Классическая траектория не выдерживает атом:
устойчивые орбиты и линейчатые спектры ей противоречат.
Опыт интерференции показывает ещё и то, что складываются не частоты исходов, а амплитуды:
иначе тёмных полос не было бы.
Состояние, из которого такие амплитуды читаются, есть вектор;
частота исхода на макроэтаже есть квадрат его модуля.
Обратимость, уже видная в уравнениях Ньютона, здесь означает,
что эволюция не меняет полную норму $\langle\psi|\psi\rangle$:
по состоянию в один момент восстанавливается состояние в предыдущий.

Опыт со спином электрона даёт два исхода и ещё одно свойство поворота.
Оборот на $2\pi$ меняет знак состояния, оборот на $4\pi$ возвращает его.
Одного комплексного числа для этого мало: минимальная запись --- столбец из двух компонент.
Общий множитель-фаза не меняет частот и соответствует сохранению заряда,
уже записанному на макроэтаже уравнением непрерывности.
Локальный поворот двух компонент лежит в $SU(2)$.

Физическая величина, построенная из такого состояния,
не имеет права зависеть от того, где проведён разрез аргумента на $2\pi$.
Комбинация, которая разреза не видит, собирается из амплитуды и сопряжённой амплитуды.
В непрерывной записи это и есть голоморфность.

Квантовая механика закрывает эти факты для одной системы с заданным числом степеней свободы.
Одиночный отсчёт прибора по-прежнему стоит рядом с уравнением.
Когда степеней свободы много и частицы рождаются и исчезают,
та же норма и тот же световой конус записываются как теория поля.

\section{Теория поля}
\label{sec:qft-open}

На масштабах ядра и элементарных частиц тот же конус, та же локальность и тот же спин
читаются как лагранжиан с калибровочными симметриями $SU(3)\times SU(2)\times U(1)$,
фермионами поколений и механизмом Хиггса.
Квантовая электродинамика и хромодинамика дают проверенные поправки.

Заряд, который доходит до прибора, кратен элементарному.
Свободный кварк не наблюдается, заряд лептона целый.
Закон, у которого устойчивый дефект нёс бы произвольную долю этого кванта,
таблицу частиц не воспроизводит.
Произведение зеркального отражения, зарядового сопряжения и обращения времени опыт сохраняет;
по отдельности отражение и зарядовое сопряжение в слабом взаимодействии нарушены.
Общий поворот фазы вакуума не меняет частот: это то же сохранение заряда.
Обе симметрии должны сидеть в самом законе.

Отсюда же видно, где непрерывная запись перестаёт справляться.

Возмущёние в квантовой теории поля содержит расходимости.
Группа ренормализации снимает их с наблюдаемых величин
и оставляет вопрос, что стоит за взаимодействием там,
где оно перестаёт быть точечным.
Это то же ограничение плотности, теперь внутри возмущения.

Стандартная модель содержит более девятнадцати настраиваемых величин:
массы, углы смешивания, константы связи.
Симметрии $SU(3)\times SU(2)\times U(1)$ их не фиксируют.
Как и этажом выше, числа стоят в законе вписанными.

Квантование метрики в той же схеме не замкнуто.
Планковская энергия~$E_P$ --- граница, за которой «поле в точке''
уже не есть запись тех же требований.
Теория поля --- последний непрерывный этаж.

\section{Планковский мир}
\label{sec:planck-floor}
\label{ch:carrier-object}
\label{sec:carrier-object}

Комбинация $\hbar$, $G$ и~$c$ задаёт единственный масштаб длины
\[
  \ell_P=\sqrt{\frac{\hbar G}{c^3}},
\]
и вместе с ним время~$t_P=\ell_P/c$ и массу~$m_P=\sqrt{\hbar c/G}$.
Это нижняя граница, на которой одновременно видны
квантовость, локальная кривизна и причинность.
Физика по-прежнему не зависит от человеческих единиц:
важны безразмерные отношения и те же симметрии, что названы выше.

Требование ограниченной плотности уже названо.
Если в каждой точке $x\in\mathbb{R}^3$ сидит своё комплексное число,
в любом конечном причинном объёме конфигураций несчётно много,
а требование уже ограничивает это число.
Место, на котором это выполнено, --- счётное множество узлов.

\begin{definition}[Носитель как множество узлов]
\label{def:carrier-sites}
Носитель микроуровня~$M$ --- счётное множество узлов~$\Lambda$.
Время, согласное с одним шагом сигнала, дискретно: $t\in\mathbb{Z}$.
Граф, рёбра которого соединяют узлы, сюда ещё не входит:
его дадут те же требования, когда они будут записаны как условия на локальный закон.
\end{definition}

Требование конечной скорости уже названо.
На счётном множестве узлов конечное расстояние за конечное время --- это один шаг за один такт.
Вторая координационная оболочка за тот же такт была бы скоростью выше этой.

\begin{definition}[Элементарные шаг и такт]
\label{def:step-tick}
$\hl$ --- элементарный шаг: длина, на которую за один такт уходит сигнал.
$\htick$ --- один такт эволюции.
Тактовая скорость сигнала на~$M$
\[
  c_0 = \frac{\hl}{\htick}.
\]
Макроскопическая скорость света~$c$ с~$c_0$ здесь не отождествляется:
безразмерный множитель между ними появится вместе с геометрией однотактового тела.
Отождествление $\hl$ с планковской длиной --- позже, когда тот же список дойдёт до системы единиц.
\end{definition}

\begin{definition}[Окрестность одного такта]
\label{def:neighborhood}
Окрестность узла $x\in\Lambda$ --- множество узлов, достижимых из~$x$
ровно за один такт по ребру длины~$\hl$:
\[
  N(x)\subset\Lambda.
\]
Число $|N(x)|$ и тип графа это определение не задаёт.
\end{definition}

Требования спина и сохранения нормы уже записаны в квантовой механике.
На узле они остаются теми же:
значение --- спинор из двух компонент, норма --- сумма квадратов модулей.

\begin{definition}[Поле и норма]
\label{def:field-a3}
Элементарная ячейка~$\PlanckCell$ --- один узел~$\Lambda$
(масштаб длины~$\hl$).
На каждом~$x$ задан спинор
\[
  z(x,t)=(z_1,z_2)^\top\in\mathbb{C}^2.
\]
Модуль $|z_1|^2+|z_2|^2$ --- плотность поля на ячейке;
фазы $z_k/|z_k|$ --- степени свободы, которые перемешиваются при эволюции.
Конфигурация всех узлов в такт~$t$ есть
\[
  \Psi^t=\{z(x,t)\}_{x\in\Lambda}.
\]
Суммарная норма, аналог $\langle\psi|\psi\rangle$,
\[
  N(t):=\sum_{x\in\Lambda} z^\dagger(x,t)\,z(x,t)=\sum_x\bigl(|z_1|^2+|z_2|^2\bigr).
\]
Сохраняет ли её шаг во времени, решает то же требование унитарности,
записанное следующей главой как условие на локальный закон.
\end{definition}

Этим спуск заканчивается.
Узлы, шаг, такт, окрестность одного такта и поле на узле
появились как то, что остаётся от уже названных требований,
когда их перестало выполнять поле в каждой точке.
Макромир, с которого начался спуск, далее обозначается~$T$:
это осреднение по~$M$, а не отдельный набор принципов.
Граф, число соседей и явный вид закона здесь ещё не написаны.
Их дадут те же требования, когда каждое будет записано как условие на допустимый шаг.

\section{Дальше по той же лестнице}
\label{sec:book-roadmap}

Следующая глава записывает то, что на этом спуске оказалось обязательным,
как условия на локальный закон в планковском мире.
Геометрия графа, запрет предела $\hl\to 0$ и явный закон эволюции
выводятся из этих условий.
Осреднение возвращает макромир, с которого спуск начался:
свободную энергию, волновое уравнение, механику ячейки
и мост к числам SI и Стандартной модели.
На обратном пути обязательство то же: осреднение обязано вернуть макромир.
